\section{Introduction}
\label{sec:introduction}

Distributing reliable entanglement across noisy quantum networks
is essential for secure communication, distributed quantum
computing, and networked quantum sensing
\cite{Wehner2018QuantumInternet}.
Shared entanglement enables quantum information processing
between distant nodes, but its usefulness is limited by noise
introduced during distribution.
Entanglement distillation and quantum network tomography (QNT)
address two complementary aspects of this problem.

Entanglement distillation uses local operations and classical
communication (LOCC) to probabilistically
 convert multiple noisy entangled pairs
into fewer pairs of higher fidelity   
\cite{PhysRevLett.76.722,Deutsch_1996}.
It provides a means to recover useful entanglement from imperfect
physical links and is an important component of quantum repeater
architectures \cite{Briegel1998Repeaters}.
For distributed quantum computing, improving shared entanglement
supports more reliable teleportation and quantum operations
between remote processors \cite{Wehner2018QuantumInternet}.

QNT infers the noise parameters of internal quantum links from
measurements performed at accessible peripheral nodes, termed
\emph{monitors} \cite{Andrade2024QNT}.
For experimental quantum networks, these estimates help diagnose
noisy links and guide distillation protocols.
For distributed quantum computing, they provide information for
predicting communication errors and selecting entanglement
distribution paths. The fidelity of an end-to-end state alone does not identify
which internal links introduce the noise; QNT seeks to recover
this link-level information from endpoint observations.
Together, entanglement distillation and QNT support
both the characterization and improvement of distributed quantum resources. Integrating them is particularly attractive when
generating entanglement is costly and measurement access is
restricted to peripheral nodes.

Despite their different objectives, entanglement distillation
and the QNT protocols considered here share a common operational feature: both rely on local measurements at accessible
nodes and classical processing of the outcomes. In distillation,
these outcomes determine whether an output pair is retained; in
QNT, their statistics are used to infer internal link parameters.
This operational similarity raises the possibility of using
distillation measurement statistics for network tomography









Casapao et al.\ pioneered the integration of entanglement
distillation and quantum state estimation with
Disti-Mator which
uses measurement statistics collected during entanglement
distillation to estimate the $X$-, $Y$-, and $Z$-basis parameters
that fully characterize a Bell-diagonal state under a fidelity
condition \cite{Casapao2025DistiMator}.
Their subsequent work developed a double-selection-based
estimator in \S~\ref{subsec:doub} using the marginal and joint coincidence statistics
collected during distillation \cite{Casapao_2025};
Motivated by these results, we identify a natural starting
point for QNT: the states shared by peripheral monitors
encode information about the individual links that we
use to distribute them.
We can therefore not only distill these states but also
estimate their parameters, and then use these estimates
to infer individual link parameters.
We thus integrate QNT into end-to-end distillation by using
distillation measurement statistics for link characterization
while also generating higher-fidelity entangled pairs,
thereby obtaining the best of both worlds.
We develop a Joint-Coincidence Estimator (JCE) and a
Marginal-Coincidence Estimator (MCE) for this task.




Realizing this idea requires overcoming two challenges. First,
state estimation does not automatically imply link identifiability.
In a three-link star network with Pauli channels, direct Unicast
transmission yields pairwise products of link-level Pauli transfer
matrix (PTM) eigenvalues. From these products, one can estimate the absolute value of the link parameters but with a sign ambiguity. To resolve this ambiguity, we adapt the Mergecast
and channel-dressing techniques introduced
in \cite{wang2025quantumnetworktomographygeneral} to prepare
Bell-diagonal states for double-selection measurements.
Mergecast provides estimates of three-link correlations,
while Clifford dressing makes the required Pauli components
accessible. Combining these estimates with those of the pairwise
products enables estimation of both the magnitudes and
signs of the link parameters without access to measurement
outcomes from the central node.

Second, inherited from the double-selection protocol, estimating all three state parameters \(q_X, q_Y, q_Z\) requires a source-state fidelity condition, implying another sign ambiguity of the protocol itself.
 This condition restricts
the noise regimes in which JCE applies. We overcome this
restriction in MCE  by observing that
the $Z$-basis parameter can be uniquely recovered without
a fidelity condition. Applying uniform Clifford dressings
maps the original $X$ and $Y$ components to the effective
$Z$ component. Repeating the estimation in these configurations
therefore enables the estimation of all link parameters
without the source-state fidelity requirement.

Our contributions are summarized as follows:
\begin{itemize}
    \item In \S~\ref{sec:dist-qnt}, we develop a distillation-assisted QNT framework for a
    three-link star network with unknown Pauli channels. The framework
    integrates Unicast and adapted Mergecast preparations with
    double-selection measurements at peripheral monitors, using
    measurement statistics from both accepted and rejected attempts for estimation.
    The central node performs prescribed operations without supplying
    measurement outcomes to the tomography procedure.

\item In \S~\ref{linkpa}, we construct two estimators to identify
the parameters of each Pauli channel in the three-link star network.
JCE combines estimates of pairwise and three-link correlations
with link-dependent Clifford dressing under the specified
source-state fidelity condition $F>1/2$.
MCE combines $Z$-basis coincidence statistics with uniform
Clifford dressing and requires no source-state fidelity condition.
\textcolor{blue}{
\item In \S~\ref{sec:numerical-experiments}, we evaluate JCE and MCE through
Qiskit simulations and compare their estimation accuracy
with existing QNT methods under different measurement
budgets and prior assumptions. The results demonstrate
the benefit of joint coincidence statistics for estimation
and show that both estimators can recover channel parameters
while retaining entangled pairs with improved fidelity.
We further extend JCE to a six-link triangular network,
demonstrating its applicability beyond the three-link
star topology.
```}

\end{itemize}

% \subsection{Related Work}
% Entanglement distillation has developed from foundational two-party
% protocols into a key tool for long-distance quantum communication.
% In 1996, Bennett et al.\ established a distillation protocol for
% extracting high-quality entangled pairs from noisy shared states
% \cite{PhysRevLett.76.722}, while Deutsch et al.\ introduced a
% purification-based quantum privacy amplification protocol
% \cite{Deutsch_1996}. Briegel et al.\ subsequently incorporated nested
% purification into quantum repeaters, connecting local entanglement
% processing with entanglement distribution over long distances
% \cite{Briegel1998Repeaters}. Further developments addressed the
% effects of imperfect local operations. In particular, Fujii and
% Yamamoto introduced double selection, which uses two ancillary pairs
% to improve error detection and achieve higher purification fidelities
% and noise thresholds than single selection in the regimes they
% studied \cite{PhysRevA.80.042308}.

% The development of QNT has focused on recovering internal link
% properties despite limited measurement access. In 2022, de Andrade
% et al.\ formulated end-to-end channel characterization through
% multipartite state distribution and studied star networks with
% single-Pauli noise, including bit-flip channels
% \cite{Andrade2022MultipartiteQNT}. Subsequent work developed
% protocols for uniquely characterizing quantum flip stars and used
% quantum Fisher information to compare their estimation performance
% \cite{Andrade2023FlipStars}. The framework was also applied to
% depolarizing channels, extending its scope beyond bit-flip noise
% \cite{Andrade2024QNT}. More recently, Wang et al.\ introduced
% Mergecast and a progressive etching procedure to address
% identifiability in networks with general Pauli channels and broader
% topologies, and further considered state-preparation and measurement
% (SPAM) errors \cite{wang2025quantumnetworktomographygeneral}.
% These studies establish methods for extracting link information
% from network-level observations, motivating the use of other
% network operations as sources of tomographic data.